\section{Introduction}

Emerging economies now account for a large share of the world's patenting: China's patent office alone received
49\% of the 3.7 million patent applications filed worldwide in 2024, up from 35\% in 2014 \citep{WIPO2025}. Yet
their patents are widely thought to be worth less. Renewal decisions offer a direct test. Owners pay a fee every year to keep a patent in force and stop paying
once it is no longer worth the fee, so the age at which patents lapse reveals what they are worth to their owners
\citep{PakesSchankerman1984,Pakes1986,SchankermanPakes1986}. By this measure the gap looks large. Among patents
filed in 2000--2005, domestic owners' patents fall short of foreign owners' patents in survival eight years after
grant by 12 percentage points more in the emerging-market offices than in the advanced ones.

Two very different stories fit such a gap. In the first, the \emph{invention} is the problem: emerging-market
patents protect weaker technology, and owners rationally drop them sooner. In the second, the \emph{inventor} is
the problem: owners lack the finance, production capacity or market access to turn valuable inventions into
returns, so good patents lapse for want of an owner able to use them. The policy responses differ. The first calls
for investment in research capability and links to frontier technology; the second for credit, fee relief or
markets in which patents can find better owners.

This paper separates the two using the legal-event records in PATSTAT,\footnote{PATSTAT is the European Patent Office's worldwide patent statistical database.} which track renewals and lapses at the
offices that report them. I first establish which offices do: of the large emerging economies, only China, Korea,
Russia, Brazil and the Czech Republic record lapses consistently and have enough domestic firm
patents.\footnote{Most countries in Latin America, Africa and South-East Asia report no lapses
to PATSTAT; Hungary reports them but has too few domestic firm patents, and Taiwan was already a high-income economy (Supplemental Appendix~S1).}
I compare these against five advanced economies (the United States, Japan, Germany, the United Kingdom, and France) using 2.5 million granted patents held by firms and universities, and renewal-fee schedules from official sources. The design compares domestic and foreign owners in the same office (exposed to the same fees, law and
enforcement), and asks whether the domestic shortfall is larger in emerging-market offices. The paper makes four
points.

First, the raw shortfall reflects what is patented. Domestic patents in emerging-market offices are cited far less
than foreign ones and are granted sooner. Comparing patents with the same citations, technology field and grant lag,
emerging-market firms keep their patents about as long as comparable firms in the advanced economies
(Section~\ref{sec:survival}): small firms are level, and mid-sized and large firms 5--8 percentage points ahead.
Office by office, Chinese small firms are level, Korean firms ahead and Russian firms
behind.

Second, three tests look for owner constraints (Section~\ref{sec:constraints}). Two find none: owners keep the home
patent of an invention longer than its foreign patents, more so in the emerging economies, and credit crunches do not
raise emerging-market firms' lapses. The third points to owner capacity at the margin: patents that small Chinese
firms sell to unrelated buyers survive 11 points longer afterwards, against little in Germany and France, although the
comparison rests on few sales there. Small Chinese firms also barely respond to fee cuts.

Third, renewal models turn these decisions into money (Section~\ref{sec:values}). Valuing each patent from its own
renewal history, without reference to its owner, the home premium (the value of domestic relative to foreign owners'
patents in the same office) does not differ significantly between the emerging-market and the advanced offices for
all owners and for large firms, and is larger in the emerging-market offices for mid-sized firms. For small firms it
is smaller against the advanced-economy pool, which Japan dominates; office by office, it is 29\% in China and 65\%
in Korea, against 28\% in the United States and 31\% in Germany.

Fourth, Chinese universities are the clear weak spot. Their patents are 15 points less likely than foreign
universities' patents in China to be alive eight years after grant and are worth 59\% less at equal citations. Here
the owner matters: once sold to an unrelated firm, Chinese university patents survive 45 points longer.

I close by pricing the policy levers on a common scale (Section~\ref{sec:stake}). Had every eligible owner taken up
China's, Korea's and Brazil's fee discounts, the returns from a cohort of patents would have risen by about a million
dollars or less. Had small Chinese firms' patents earned what they earn in buyers' hands, returns would have risen by up to
\$130 million per cohort, or \$280 million with mid-sized firms; moving Chinese university patents to firms is worth
up to \$460 million, and raising their returns to what foreign universities' patents earn in China at equal
citations, \$180--320 million. These are
rough private returns, not welfare, but the ordering is clear: the gains lie in markets for technology and in what
universities patent, not in cheaper renewal.

The paper contributes to three literatures. The first estimates the value of patent rights from renewal decisions
\citep{PakesSchankerman1984,Pakes1986,SchankermanPakes1986,Lanjouw1998,Schankerman1998,LanjouwPakesPutnam1998,Bessen2008}
and the role of fees \citep{deRassenfosseVanPottelsberghe2013}; I extend it to emerging-market
offices, with fee histories for ten offices over 25 years. The second studies indicators of patent value
\citep{HarhoffSchererVopel2003,HallJaffeTrajtenberg2005,KoganEtAl2017}; I value each patent from its own renewal
path rather than from its owner's characteristics. The third asks whether Chinese and other emerging-market patenting
is of low quality \citep{HuJefferson2009,LiegsalzWagner2013,DangMotohashi2015,HuZhangZhao2017,WeiXieZhang2017}; the answer here is that, for firms,
lower survival on average reflects the technology rather than the owners' ability to use their patents, although
owner capacity matters at the margin. The sale premium connects to
work on markets for technology
\citep{AroraFosfuriGambardella2001,Serrano2010,Serrano2018,GalassoSchankermanSerrano2013,AkcigitCelikGreenwood2016},
and the university result to work on who should own university inventions \citep{HvideJones2018}.
